\chapter{Coughing Up Phlegm}

\section*{Definition}
Coughing up phlegm is productive cough that brings mucus from the airways. Color alone does not reliably distinguish viral from bacterial infection.

\section*{When to See a Doctor}
Seek urgent evaluation for Breathlessness, chest pain, coughing blood, blue lips, confusion, high fever, dehydration, or low oxygen. Arrange routine or prompt assessment for persistent, recurrent, unexplained, or function-limiting symptoms.

\section*{How It Develops}
Mucus is produced to protect the airways and is normally cleared by cilia and swallowing. Infection, airway inflammation, smoking, chronic lung disease, or impaired clearance can increase mucus and make coughing bring it up.

\section*{Causes}
\begin{itemize}
  \item Viral bronchitis, pneumonia, asthma, COPD, bronchiectasis, postnasal drainage, smoking, and chronic infection.
  \item Medication effects, environmental exposures, and chronic disease may contribute.
  \item Less common but important causes include malignancy, vascular disease, or organ failure when symptoms persist or progress.
\end{itemize}

\section*{Related Symptoms}
\begin{itemize}
  \item Wheezing, breathlessness, fever, chest discomfort, nasal drainage, or fatigue
  \item Foul-smelling or persistent sputum, weight loss, night sweats, or coughing blood
\end{itemize}

\section*{Medical Evaluation}
\begin{itemize}
  \item History addresses onset, duration, severity, triggers, exposures, medications, medical conditions, pregnancy status when relevant, and warning symptoms.
  \item Examination includes vital signs, oxygen saturation, and the nose, throat, and lungs.
  \item Chest radiography, spirometry, sputum testing, or other investigations are selected according to duration, severity, examination findings, and suspected cause; sputum color alone does not establish a bacterial infection.
\end{itemize}

\section*{Treatment Options}
Treatment depends on the cause. Hydration and airway-clearance measures may help; antibiotics are reserved for selected bacterial infections, and chronic lung disease requires its prescribed management. Persistent or bloody sputum needs evaluation.

\section*{Self-Care and Home Management}
Avoid known triggers, maintain reasonable hydration, record timing and associated symptoms, and use only treatments appropriate for the suspected cause. Do not use leftover antibiotics or delay urgent care because symptoms temporarily improve.

\section*{References}
\begin{thebibliography}{99}
\bibitem{coughing_up_phlegm:accp} Irwin RS, Baumann MH, Bolser DC, et al. Diagnosis and management of cough executive summary: ACCP evidence-based clinical practice guidelines. \textit{Chest}. 2006;129(suppl 1):1S--23S.
\bibitem{coughing_up_phlegm:ers} Morice AH, Millqvist E, Bieksiene K, et al. ERS guidelines on the diagnosis and treatment of chronic cough in adults and children. \textit{Eur Respir J}. 2020;55:1901136.
\bibitem{coughing_up_phlegm:harrison} Jameson JL, et al. \textit{Harrison's Principles of Internal Medicine}. 21st ed. McGraw-Hill; 2022.
\end{thebibliography}
